\song{Pacholek a děvečka (moravská lidová balada)}
%Mara_2017_02_06

\ve{}
V \ch{Emi}{š}irém poli stojí zá\ch{D}mek, v \ch{Emi}{š}irém poli \ch{D}stojí zámek\\
\ch{Emi}Slouží tam je\ch{H7}den pacholek, \ch{Emi}slouží tam je\ch{H7}den pacho\ch{Emi}lek

\begin{multicols*}{2}
\ve{}
Slouží tam s jednou děvečkou,\\
to s rychtářovou dcérečkou.

\ve{}
Sedm let tam spolu byli,\\
slova spolu nemluvili.

\ve{}
A jak přišlo léto osmé,\\
dal jim pán Bůh dítě krásné.

\ve{}
Tak se spolu domluvili,\\
že by to dítě zabili.

\ve{}
Dělej, milá, jak jen myslíš,\\
jenom ať se nevyzradíš.

\ve{}
Zlatou šňůru ukroutila,\\
do studnice ho spustila.

\ve{}
Děvče pána to vidělo,\\
hned pánovi žalovalo.

\ve{}
Pohleď Aničko do pole,\\
co uvidíš, to je tvoje.

\ve{}
Stojí, stojí šibenička,\\
na ní bílá holubička.

\ve{}
Kdo teď na ní viset bude,\\
kdo by jiný, než Anička.

\ve{}
Už ji vede kat po mostě,\\
a rozmlouvá s ní po sprostě.

\ve{}
Chceš, Aničko, moje býti,\\
já ti koupím živobytí.

\ve{}
Dělej, kate, co máš dělat,\\
moje očka zavazovat.

\columnbreak
\vspace*{\fill}

\ve{}
V širém poli stojí zámek,\\
slouží tam jeden pacholek.

\end{multicols*}

\esong